\section{The integral parameter relation in original cote 14}
\label{sec:cc-integral-albanese}

\paragraph{Scope and convention.}
Let $k$ be algebraically closed of characteristic $p>0$ and
let $T/k$ be a smooth connected quasiprojective variety.
Write
\[
 H_T=H^S_0(T,\mathbb Z)^0,\qquad
 A_T=\operatorname{Alb}_{\rm ab}(T),\qquad
 K_T=\ker\bigl(H_T\longrightarrow A_T(k)\bigr).
\]
Here $\operatorname{Alb}_{\rm ab}$ means the ordinary abelian
Albanese: after choosing $t_0\in T(k)$, its pointed morphism
$a:T\to A_T$ is universal for pointed morphisms to
\emph{abelian varieties}.  It is the abelian quotient of the
generalized semiabelian Albanese.  We use the integral Suslin
complex of finite correspondences.  In particular $H_T$ is a
quotient of the group of free degree-zero zero cycles, and
the Albanese homomorphism is the sum of their images.

The new deduction below proves divisibility of $K_T$ and
then removes the possible $p$-primary discrepancy left by
AQ-14-ALTERATION.  It requires no resolution of singularities
or smooth compactification of $T$.  It does not assert
injectivity of the Albanese map on $p$-primary torsion.

\paragraph{Exact cited inputs.}
Katz (1993), Corollary 3.4.3, printed p.~182, gives smooth
irreducible affine hyperplane sections of a smooth
irreducible variety equipped with an unramified morphism to
affine space; the sections can simultaneously preserve
irreducibility of any specified finite \'{e}tale cover.
Its hypotheses are over an algebraically closed field of
characteristic different from the auxiliary coefficient
prime.  We choose that prime different from $p$.
Suslin and Voevodsky (1996), Theorem 3.1, printed p.~65,
identifies the integral zeroth Suslin homology of a smooth
affine curve $C$ with
$\operatorname{Pic}(\overline C,\overline C\setminus C)$,
where the boundary is reduced.
The generalized Albanese universal property is recalled
in Spie{\ss} and Szamuely (2000), \S3, pp.~4--5;
passing to its abelian quotient gives the ordinary
universal property used here.
For the auxiliary curve-through-points construction we
use the Stacks Project, Lemmas 33.47.1 (Tag 0FD5) and
37.32.3 (Tag 0G4F).
Quotients by finite group schemes, their torsor property,
the connected--\'{e}tale sequence, and finiteness and
surjectivity of multiplication on abelian varieties are
recorded in Milne (2017), Theorem 5.28 and Corollaries
5.26--5.27, p.~105, Proposition 5.58, p.~114, and 8.21,
p.~152.  These statements include nonreduced subgroup
schemes; treating only abstract groups of points here
would lose the characteristic-$p$ information.

\subsection{CC-ALB01: divisible ordinary Albanese kernel}

\paragraph{Statement.}
For every positive integer $n$ the map
\[
 H_T[n]\longrightarrow A_T(k)[n]
 \quad\text{is surjective}.                         \tag{CC.ALB.1}
\]
The group $H_T$ is divisible, and its subgroup $K_T$ is
divisible.

\begin{proof}
\emph{1. Curve classes make $H_T$ divisible.}
First, any finite set of points in an integral
quasiprojective variety over $k$ lies on an integral
curve, unless the variety has dimension zero.
Here is the precise Bertini reduction used for this fact.
Take an integral projective closure $Y$, let $F$ be the
finite set, and use sections of
$\mathcal I_F\otimes\mathcal L^{\otimes m}$ for $m$
sufficiently large and $\mathcal L$ ample.
The cited Tag 0FD5 says that this linear system has no
base points outside $F$ and gives an immersion there.
The base locus on $Y$ is exactly $F$, because the ideal
sheaf is generated after a sufficiently high twist.
If $\dim Y\geq2$, this base locus has codimension at least
two.  At a smooth point outside $F$, immersion gives one
section nonzero at that point and two sections whose
local equations have independent differentials.
Their intersection has a codimension-two component
through that point, not contained in the first section.
These verify conditions (3) and (4), as well as (1) and
(2), of the cited Tag 0G4F.  Thus a general member is
geometrically irreducible.  Its reduction contains $F$
and has dimension one less.  Iterating gives an
integral curve through $F$ in the projective closure;
its nonempty intersection with the original variety
is the required integral curve.

Normalize this curve and choose one preimage of each
point of $F$.  Normalization is finite, and a normal
curve over the perfect field $k$ is smooth.  If necessary
remove a finite set disjoint from the chosen preimages
so that the resulting curve $C$ is affine.  Thus every
free degree-zero cycle on $T$ is the image of a
degree-zero cycle on some smooth affine curve $C$.
All residue fields are $k$, so the degree is unchanged.

Let $\overline C$ be the smooth projective completion,
let $D=\overline C\setminus C$ with its reduced structure,
and put $J=\operatorname{Pic}^0(\overline C)$.
The curve calculation in the cited Suslin--Voevodsky
theorem and the elementary exact sequence for line
bundles with a trivialization on $D$ give
\[
 0\longrightarrow (k^\times)^{|D|-1}
 \longrightarrow H^S_0(C,\mathbb Z)^0
 \longrightarrow J(k)\longrightarrow0.             \tag{CC.ALB.2}
\]
Indeed, a line bundle restricts to a trivial line bundle
on the finite reduced set $D$; two choices of
trivialization differ by an element of
$(k^\times)^D/k^\times$.
Both outer groups in (CC.ALB.2) are divisible:
this is immediate for $k^\times$, and multiplication
by $n$ on an abelian variety is surjective.
An extension of divisible groups is divisible: divide
the image in the quotient, lift the division, and
correct the remaining error by division in the kernel.
The functorial pushforward to $H_T$ therefore divides
each degree-zero cycle class by any prescribed $n$.
This proves divisibility of $H_T$.  If $\dim T=0$,
both $H_T$ and $A_T$ are zero and all assertions follow.
We henceforth suppose $\dim T\geq1$.

\emph{2. Albanese universality makes isogeny pullbacks
connected.}
Put $A=A_T$ and normalize $a(t_0)=0$.
Let $\pi:B\to A$ be any finite \'{e}tale isogeny of
abelian varieties.  Its kernel is the constant finite
abelian group scheme associated with a group $G$.
The pullback
\[
 P=T\times_{A,\pi}B
\]
is a finite \'{e}tale $G$-torsor.
Suppose it is disconnected.  Its connected components
are permuted transitively by $G$: each component maps
surjectively to the connected base $T$, and the action
on each geometric fibre is simply transitive.
Translate a chosen component so that it contains
$(t_0,0)$, and let $G_0\subsetneq G$ be its stabilizer.
The chosen component is then a $G_0$-torsor over $T$.
Its map to $B/G_0$ descends to a pointed morphism
\[
 b:T\longrightarrow B/G_0,\qquad \overline\pi b=a,
\]
where $\overline\pi:B/G_0\to A$ is an \'{e}tale isogeny
with nontrivial kernel $G/G_0$.

The target $B/G_0$ is an abelian variety.
The universal property used is precisely that of the
\emph{ordinary abelian} Albanese, so it supplies a
homomorphism $r:A\to B/G_0$ such that $ra=b$.
Uniqueness in that property gives
$\overline\pi r=\operatorname{id}_A$.
This is impossible for a nontrivial isogeny:
$r$ has trivial kernel and closed image of dimension
$\dim A=\dim(B/G_0)$, hence is surjective;
a point of $\ker\overline\pi$ then has the form $r(x)$,
and applying $\overline\pi$ forces $x=0$.
The nontrivial \'{e}tale group $G/G_0$ has nonzero
$k$-points, giving a contradiction.
Consequently $P$ is connected.
As $T$ and $P$ are smooth, they are also irreducible.

\emph{3. Separate the connected part of multiplication.}
Fix $n\geq1$ and let $F_n=A[n]^\circ$ be the connected
component of the identity in the finite group scheme
$A[n]$.  The cited group-scheme quotient results give
the factorization
\[
 A\xrightarrow{r_n}B_n=A/F_n
   \xrightarrow{\pi_n}A,\qquad
 \pi_n r_n=[n]_A.                                  \tag{CC.ALB.3}
\]
The quotient $B_n$ is a smooth connected proper
commutative group, hence an abelian variety.
The map $r_n$ is finite faithfully flat with geometrically
connected, single-point fibres, so it is a universal
homeomorphism.  The kernel of $\pi_n$ is the finite
\'{e}tale group scheme $A[n]/F_n$.  Thus $\pi_n$ is a
finite \'{e}tale isogeny.  These assertions also hold
when this kernel is trivial.
Step 2 shows that $T\times_A B_n$ is connected and
irreducible.

\emph{4. Keep this one cover connected on a curve.}
Choose a nonempty affine open $V\subset T$ and embed
it into affine space.  The restricted cover
$V\times_A B_n$ remains irreducible.  If $\dim V>1$,
apply Katz's Corollary 3.4.3 with $f=0$ and the closed
immersion $V\hookrightarrow\mathbb A^N$ as its
unramified morphism.  A choice in the nonempty
intersection of the two open parameter sets gives a
smooth irreducible affine hyperplane section whose
pullback cover is irreducible.  Iterate until obtaining
a smooth affine curve $C\subset V$ with
$C\times_A B_n$ connected.  If $\dim V=1$, take
$C=V$.  This is a finite sequence of choices for the
single specified $n$; no simultaneous choice for
infinitely many covers is claimed.

Let $\overline C$ be its smooth projective completion.
Properness of $A$ extends $a|_C$ over $\overline C$.
Choose $c_0\in C(k)$ and replace this map by
$a|_{\overline C}-a(c_0)$.  Translation by a lift of
$a(c_0)$ in $B_n(k)$ identifies the old and new
pullbacks, so connectedness is unchanged.
The universal property of the Jacobian gives
\[
 f:J=\operatorname{Jac}(\overline C)\longrightarrow A
\]
whose composite with $c\mapsto[c-c_0]$ is the
normalized map.  The finite \'{e}tale cover
$J\times_A B_n$ is connected as well.
Indeed, if it were disconnected, each of its
nonempty open-and-closed components would map
surjectively to the connected variety $J$.
Their pullbacks to $C$ would be nonempty disjoint
open-and-closed pieces, contradicting the preceding
connectedness.

\emph{5. Compute the components and obtain torsion
surjectivity.}
Consider
\[
 Q=J\times_{A,[n]_A}A
   =\{(x,z):f(x)=[n]_A z\}.
\]
The map $Q\to J\times_A B_n$ induced by $r_n$ is
the base change of the universal homeomorphism
$r_n$ along the projection to $B_n$.  Hence $Q$ is
connected, even though it may be nonreduced.
For each $t\in A[n](k)$, let $L_t$ be the image,
on underlying topological spaces, of the proper
morphism
\[
 J\longrightarrow Q,\qquad
 y\longmapsto\bigl([n]_J y,\ f(y)+t\bigr).           \tag{CC.ALB.4}
\]
Every $L_t$ is closed and connected.
These finitely many sets cover $Q$: given a
$k$-point $(x,z)$, choose $y$ with $[n]_J y=x$;
then $t=z-f(y)\in A[n](k)$.
This covers every closed point, hence all of the
finite-type $k$-scheme $Q$.
Moreover,
\[
 L_t=L_{t'}
 \quad\Longleftrightarrow\quad
 t-t'\in f(J[n](k)),                               \tag{CC.ALB.5}
\]
and any two intersecting $L_t,L_{t'}$ are equal.
For the latter assertion, an intersection contains
a closed point; two preimages $y,y'$ then have
$[n](y-y')=0$ and $f(y-y')=t'-t$.
Conversely translation by any such element of
$J[n](k)$ identifies the two images.
Thus the distinct $L_t$ are a finite disjoint closed
cover and are exactly the connected components.
Connectedness of $Q$ proves
\[
 f:J[n](k)\twoheadrightarrow A[n](k).               \tag{CC.ALB.6}
\]

The exact sequence (CC.ALB.2) for this curve gives a
surjection
$H^S_0(C,\mathbb Z)^0[n]\twoheadrightarrow J[n](k)$:
lift an $n$-torsion point of $J(k)$, then subtract
an $n$-division of the lift's $n$-multiple in the
torus kernel.  Functorial pushforward along $C\to T$
and (CC.ALB.6) prove (CC.ALB.1).
Translations introduced by the choice of $c_0$ vanish
on degree-zero cycles, so this is exactly the original
ordinary Albanese homomorphism.

\emph{6. Divide inside the kernel.}
Let $x\in K_T$ and let $n\geq1$.  By Step 1 choose
$y\in H_T$ with $ny=x$.
Then $\operatorname{alb}_T(y)\in A_T(k)[n]$.
By (CC.ALB.1) choose $z\in H_T[n]$ with the same
Albanese image.  The element $y-z$ belongs to $K_T$
and satisfies $n(y-z)=x$.  This proves divisibility
of $K_T$.
\end{proof}

\subsection{CC-ALB02: integral equality of projective and
quasiprojective parameter relations}

\paragraph{Statement.}
Let $X/k$ be smooth projective over an algebraically
closed field, of arbitrary characteristic.
For a fixed codimension $i$, let
$N^i_{\rm proj}(X)$ be generated in $CH^i(X)$ by
\[
 \Gamma_*v,\qquad
 \Gamma\in CH^i(X\times T),\quad
 v\in Z_0(T)^0,\quad \operatorname{alb}_T(v)=0,
                                                               \tag{CC.ALB.7}
\]
with $T$ smooth connected projective.
Define $N^i_{\rm qp}(X)$ in the same way with smooth
connected quasiprojective $T$, using its ordinary
abelian Albanese.  Then
\[
 N^i_{\rm proj}(X)=N^i_{\rm qp}(X).                 \tag{CC.ALB.8}
\]
The same subgroup results if the quasiprojective
definition uses the generalized semiabelian
Albanese kernel.

\begin{proof}
In characteristic zero this is the already proved
smooth compactification argument of
AQ-14-QUASIPROJECTIVE.  Suppose $\operatorname{char}k=p>0$.
Fix one parameter $T$ and one correspondence $\Gamma$.
The proper-target Suslin descent proved in
AQ-14-NONPROPER makes its action factor through
$H_T$.  We therefore have a homomorphism
\[
 \varphi_\Gamma:K_T\longrightarrow
 CH^i(X)/N^i_{\rm proj}(X).
\]

We make explicit the uniform quantifier in the
already proved alteration argument.
Choose once a projective compactification
$T\subset\overline T$ and a separable projective
alteration $S\to\overline T$ with $S$ smooth
projective, of degree $d>0$.
For example, choose any auxiliary prime $\ell\ne p$
and apply Illusie--Temkin (2014), Expos\'e X,
Theorem 3.5(iii), p.~202, as in AQ-14-ALTERATION.
There is a single nonempty open $U\subset T$ over
which the alteration is finite \'{e}tale of degree $d$.

For every $x\in K_T$, the moving proof in Step 1 of
AQ-14-ALTERATION gives a representative free
degree-zero cycle $v_U$ supported in this same $U$,
with Albanese sum zero and the same Suslin class.
The fibre-sum morphism
$U\to\operatorname{Alb}(S)$ extends across smooth
$T$ and factors, up to translation, through
$\operatorname{Alb}_{\rm ab}(T)$.
Consequently the pullback cycle $w=f_U^*v_U$ has
Albanese sum zero on $S$.
Closing the pullback of $\Gamma|_{X\times U}$
on $X\times S$ gives a correspondence $\Gamma'$
and the fibre Gysin calculation gives
\[
 \Gamma'_*w=d\,\Gamma_*v_U.
\]
The left side lies in $N^i_{\rm proj}(X)$.
Thus the \emph{same} positive integer $d$, fixed
before choosing $x$, satisfies
\[
 d\,\varphi_\Gamma(x)=0\qquad\text{for every }x\in K_T.
                                                               \tag{CC.ALB.9}
\]
Only the moved cycle varies with $x$; the alteration
and its degree do not.

By CC-ALB01, write $x=dy$ with $y\in K_T$.
Equation (CC.ALB.9) yields
$\varphi_\Gamma(x)=d\varphi_\Gamma(y)=0$.
This holds for every $T,\Gamma,x$, so every generator
of $N^i_{\rm qp}(X)$ belongs to $N^i_{\rm proj}(X)$.
The reverse inclusion is immediate, proving
(CC.ALB.8).

Finally the generalized semiabelian Albanese maps
to its ordinary abelian quotient.  Its kernel
relation is therefore contained in the relation
using the ordinary Albanese.  On projective
parameters the generalized Albanese is already
abelian.  The semiabelian parameter subgroup lies
between the two equal groups in (CC.ALB.8), so it
is equal to them.
\end{proof}

\paragraph{What the deduction settles.}
For the free-zero-cycle parameter convention actually
used in the cited cote 14 formulation, the possible
integral $p$-primary discrepancy between smooth
projective and smooth quasiprojective parameters
vanishes.  The earlier away-from-$p$ statement remains
a valid intermediate result.  No statement about
injectivity of the Suslin Albanese map on
$p$-primary torsion, integral motivic projectors,
or resolution of singularities is needed or inferred.

\paragraph{References and inspected locators.}
\begin{itemize}
\item Katz, N. M. (1993). Affine cohomological
transforms, perversity, and monodromy.
\emph{Journal of the American Mathematical Society,
6}(1), 149--222. Corollaries 3.4.2--3.4.3,
pp.~181--182.
\url{https://web.math.princeton.edu/~nmk/old/ACT.pdf}
\item Suslin, A., \& Voevodsky, V. (1996).
Singular homology of abstract algebraic varieties.
\emph{Inventiones Mathematicae, 123}, 61--94.
Theorem 3.1, p.~65, proof pp.~65--66.
\url{https://doi.org/10.1007/BF01232367}
\url{https://www.math.ias.edu/vladimir/sites/math.ias.edu.vladimir/files/singular_homology_published.pdf}
\item Milne, J. S. (2017). \emph{Algebraic groups:
The theory of group schemes of finite type over a
field}. Cambridge University Press.
Theorem 5.28 and Corollaries 5.26--5.27, p.~105;
Proposition 5.58, p.~114; 8.18 and 8.21, p.~152.
\url{https://www.jmilne.org/math/Books/iAG2017.pdf}
\item The Stacks Project Authors. (2026).
\emph{The Stacks Project}. Lemma 33.47.1,
Tag 0FD5, and Lemma 37.32.3, Tag 0G4F.
\url{https://stacks.math.columbia.edu/tag/0FD5}
\url{https://stacks.math.columbia.edu/tag/0G4F}
\item Spie{\ss}, M., \& Szamuely, T. (2000).
\emph{On the Albanese map for smooth quasi-projective
varieties} (arXiv:math/0009017v1). arXiv.
\S3, pp.~4--5, generalized Albanese universal property.
\url{https://arxiv.org/pdf/math/0009017v1}
\item Illusie, L., \& Temkin, M. (2014).
Expos\'e X: Gabber's modification theorem
(log smooth case).
In L. Illusie, Y. Laszlo, \& F. Orgogozo (Eds.),
\emph{Travaux de Gabber sur l'uniformisation locale
et la cohomologie \'{e}tale des sch\'emas
quasi-excellents} (Ast\'erisque 363--364,
pp.~167--212). Soci\'et\'e Math\'ematique de France.
Theorem 3.5(iii), p.~202; the application and
separability were already inspected in
AQ-14-ALTERATION.
\url{https://www.numdam.org/item/AST_2014__363-364__167_0.pdf}
\end{itemize}
